% \section{Quasi-Hyperbolic Discounting}
\label{app:QuasiHyperbolic}
Quasi-hyperbolic preferences are a tractable way to model 'impatience', with the present-self taking decisions that are not in the interest of their own future-self. The standard way to discount the future is exponential discounting, which uses the same discount factor $\beta$ between any two consecutive periods. Quasi-hyperbolic discounting involves two discount factors, $\beta_0$ which is used as the additional discount factor between the current period and the next period, and $\beta$ which is used as the discount factor between any two consecutive periods; so the discount factor between the current period and the next period is $\beta_0 \beta$.\footnote{We use $\beta_0 \beta$ as the discount rate between this period and next period, rather than just some $\alpha \equiv \beta_0 \beta$ because this notation becomes much more convenient when we want to look at, e.g., life-cycle models where there is a probability of dying; especially in terms of writing the codes. Note that papers that look at theory on Quasi-Hyperbolic discounting in simple models often just use one parameter to control this period to next period, e.g., $\alpha \equiv \beta_0 \beta$.} The concept of hyperbolic discounting has a long heritage, and the tractable case of quasi-hyperbolic was popularised (not introduced) in \cite*{Laibson1997}. Let's introduce the idea using a four period model. I first will describe quasi-hyperbolic here in a deterministic setting (without uncertainty), and then switch to a stochastic setting when switching to recursive notation

In a four period model with standard exponential discounting
\begin{equation}
 u(c_1)+\beta u(c_2) +\beta^2 u(c_3)+\beta^3 u(c_4) \label{eqn:expdiscount}
\end{equation}
Notice how $\beta$ is used as the discount rate between periods one-and-two as well as between periods two-and-three and three-and-four. In contrast quasi-hyperbolic discounting uses $\beta_0 \beta$ between periods one-and-two, but then uses $\beta$ between periods two-and-three and three-and-four. Thus our four period model with quasi-hyperbolic discounting becomes
\begin{equation}
 u(c_1)+\beta_0 \beta u(c_2) +\beta_0  \beta^2 u(c_3)+ \beta_0 \beta^3 u(c_4) \label{eqn:qhdiscount}
\end{equation}
Because we have not fully specified the optimization problem (what is being choosen, just $c_1$? or also $c_2$, $c_3$, \& $c_4$?) it is not so obvious in this formulation that there are two types of quasi-hyperbolic discounter: naive and sophisticated. A \textit{naive} quasi-hyperbolic discounter assumes that their future self will not be impatient even though their present self is (their future self is using exponential discounting). A \textit{sophisticated} quasi-hyperbolic discounter; one who realises that their future self will also behave impatiently by discounting in a quasi-hyperbolic manner. This will be much clearer when we now switch to recursive notation and write out the optimization problem in full.

We now express the same thing in recursive notation, but adding stochastics ($z$ and the expectations). We have that standard exponential discounting can be expressed as the value function
\begin{equation} \label{eqn:naive1}
    V(a,z)=max_{c,a' \in D(a,z)} u(c)+\beta E[V(a',z')]
\end{equation}
where $a$ is an endogenous state, and $z$ is an exogenous stochastic state. $c$ and $a'$ are constrainted to be in some feasible decision space $D(a,z)$. In finite time $V(a',z')$ would be next period value function and so not the same as $V(a,z)$, in infinite time they are the same.

We can express naive quasi-hyperbolic discounting in terms of the exponential discounting solution. Let $\tilde{V}$ be the value function of the naive quasi-hyperbolic discounter, then
\begin{equation} \label{eqn:naive2}
    \tilde{V}(a,z)=max_{c,a' \in D(a,z)} u(c)+\beta_0 \beta E[V(a',z')]
\end{equation}
notice how the expected next-period value function is that of the behavior of the exponential discounter, who the naive quasi-hyperbolic discounter (incorrectly) believes they will act like. We call $V$ the 'continuation value' of the naive quasi-hyperbolic discounter.

The sophisticated quasi-hyperbolic discounter understands that their future self will behave as a quasi-hyperbolic discounter. Let $\hat{V}$ be the value function of the sophisticated quasi-hyperbolic discounter, then
\begin{equation} \label{eqn:sophisticated1}
    \hat{V}(a,z)=max_{c,a' \in D(a,z)} u(c)+\beta_0 \beta E[\underbar{V}(a',z')]
\end{equation}
and the argmax of this same expression gives the (optimal) policies $\hat{c}$ and $\hat{a}'$. Using $\hat{c}$ and $\hat{a}'$ we get the definition of $\underbar{V}(a,z)$ as
\begin{equation} \label{eqn:sophisticated2}
    \underbar{V}(a,z)= u(\hat{c})+\beta E[\underbar{V}(\hat{a}',z')]
\end{equation}
note that to compute this we will need, given next period $\underbar{V}$, to first get current period $\hat{V}$ and $\hat{c}$ and $\hat{a}'$ from equation (\ref{eqn:sophisticated1}), and then use these in equation (\ref{eqn:sophisticated2}) to calculate the current period $\underbar{V}$. We call $\underbar{V}$ the 'continuation value' of the naive quasi-hyperbolic discounter.\footnote{One way to see is going on with the continuation value of the sophisticated quasi-hyperbolic discounter is to return our trivial example of $u(c_1)+\beta_0 \beta u(c_2) +\beta_0  \beta^2 u(c_3)+ \beta_0 \beta^3 u(c_4)$, and rewrite this as $u(c_1)+\beta_0 \beta [u(c_2) +\beta u(c_3)+ \beta u(c_4)]$, now we can see that we have the current problem $u(c_1)+\beta_0 \beta 'continuation \, value'$, and the continuation value is $u(c_2) +\beta u(c_3)+ \beta u(c_4)$  which only uses $\beta$ and for the sophisticated agent they know that this continuation value, their $underbar{V}$, is based on the actions that an impatient agent will take (the 'hat' policies).}

Notice that standard exponential discounting is nested as the case $\beta_0=1$, and quasi-hyperbolic discounting involves $0<\beta_0<1$. For infinite horizon models $0<\beta<1$, but for finite horizon models with a probability of dying it can sometimes be greater than 1.

Quasi-hyperbolic preferences are time-inconsistent preferences.
